\documentclass[11pt,a4paper]{article}
\usepackage[T1]{fontenc}
\usepackage[utf8]{inputenc}
\usepackage{amsmath,amssymb,amsthm}
\usepackage[margin=1in]{geometry}
\usepackage[colorlinks=true,linkcolor=blue,urlcolor=blue]{hyperref}
\theoremstyle{plain}
\newtheorem{theorem}{Theorem}
\newtheorem{lemma}[theorem]{Lemma}
\newtheorem{proposition}[theorem]{Proposition}
\newtheorem{corollary}[theorem]{Corollary}
\theoremstyle{definition}
\newtheorem{remark}[theorem]{Remark}

\title{The empirical recurrence for OEIS A184021, proved by splitting on the common sum}
\author{Adrian Perez Fontelles\\ \small Independent researcher}
\date{1 September 2026}
\begin{document}
\maketitle

\begin{abstract}
OEIS A184021 counts arrays over $\{0,\dots,3\}$ in which all $2\times2$ subblock sums are
EQUAL --- the common value not being named --- and carries an empirical recurrence of order
$15$ contributed by R.~H.~Hardin. The unnamed value is what stops a transfer matrix from
applying directly, and it costs nothing to remove: an array with at least one $2\times2$
subblock determines its own common sum, so the count splits as a sum over the $12+1$
possible values with nothing counted twice, and each summand is an ordinary walk count. On
the block-diagonal matrix over those values the recurrence is decided by a finite exact
computation on $S=208$ vertices.
\end{abstract}

\noindent\small 2020 Mathematics Subject Classification. 05A15, 05B45, 15B36.\normalsize

\section{The sequence and the conjecture}

OEIS A184021 is ``1/8 the number of (n+1)X2 0..3 arrays with all 2X2 subblock sums the same.''. It has offset $1$ and begins
\[
32, 88, 242, 748, 2312, 7684, 25538, 89044,\ \dots
\]
The entry carries the following comment:
\begin{quote}\small
Empirical: a(n)=8*a(n-1)+16*a(n-2)-268*a(n-3)+189*a(n-4)+3388*a(n-5)-5824*a(n-6)-20048*a(n-7)+49756*a(n-8)+52752*a(n-9)-196896*a(n-10)-24192*a(n-11)+369216*a(n-12)-131328*a(n-13)-262656*a(n-14)+165888*a(n-15)
\end{quote}
As of the ``Last modified'' line on the live entry (Oct 03 2025 02:15:02, revision 9) the statement is
still recorded as empirical, and nothing on the entry records it as settled either way.

Written out, the claim is
\begin{equation}\label{eq:conj}
a(n)\;=\;8\,a(n-1) + 16\,a(n-2) - 268\,a(n-3) + 189\,a(n-4) + 3388\,a(n-5) - 5824\,a(n-6) - 20048\,a(n-7) + 49756\,a(n-8) + 52752\,a(n-9) - 196896\,a(n-10) - 24192\,a(n-11) + 369216\,a(n-12) - 131328\,a(n-13) - 262656\,a(n-14) + 165888\,a(n-15)
\end{equation}
for all $n>15$.

\section{Splitting on the common sum}

The arrays have $2$ columns and $L=n+1$ rows, entries in
$\{0,\dots,3\}$, and every $2\times2$ subblock has the same sum. Since $L\ge2$ and
$2\ge2$, every such array contains at least one $2\times2$ subblock, and that subblock's
sum is the common value; so each array belongs to exactly one of the classes
\[
A_c=\{\text{arrays all of whose }2\times2\text{ subblock sums equal }c\},
\qquad 0\le c\le 12,
\]
and $\#\{\text{admissible arrays}\}=\sum_c |A_c|$ with no array counted twice. The bound
$12=4\cdot3$ is the largest possible subblock sum.

For a fixed $c$ the condition is local, so $|A_c|$ is a walk count. Let
$V=\{0,\dots,3\}^{2}$ be the possible rows and let $M_c$ be the adjacency
matrix of the digraph in which $r\to s$ when
$r_j+r_{j+1}+s_j+s_{j+1}=c$ for every $j$ with $1\le j\le 2-1$. Then
$|A_c|=\mathbf{1}^{\!\top}M_c^{\,L-1}\mathbf{1}$, since an array of $L$ rows is a
sequence of $L$ elements of $V$ and each of its $2\times2$ subblocks lies in two consecutive
rows and two consecutive columns.

Writing $N=\bigoplus_{c=0}^{12}M_c$, of size $S=208$, and $\mathbf{v}$ for the
all-ones vector on it,
\begin{equation}\label{eq:walk}
a(n)\;=\;\tfrac1{8}\,\mathbf{v}^{\!\top}N^{\,L-1}\mathbf{1},\qquad L=n+1.
\end{equation}

\section{The criterion, and the computation}

Let $q(t)=t^{15}-\sum_i c_i\,t^{15-i}$ be the characteristic polynomial of
\eqref{eq:conj} and $G=N$.

\begin{lemma}\label{lem:crit}
With $n_0$ the least index for which $L\ge1$, $o=1n_0+1-1$ and
$h_j=G^{j}N^{o}\mathbf{1}$, we have for $j\ge15$
\[
a(n_0+j)-\sum_i c_i\,a(n_0+j-i)=\tfrac1{8}\,\mathbf{v}^{\!\top}G^{\,j-15}w,
\qquad w=h_{15}-\sum_i c_ih_{15-i} .
\]
Hence \eqref{eq:conj} holds from that point on if and only if
$\mathbf{v}^{\!\top}G^{m}w=0$ for all $m\ge0$, and $m<S$ suffices.
\end{lemma}

\begin{proof}
Substitute \eqref{eq:walk} and factor. By Cayley--Hamilton $G^{S}$ is an integer
combination of $I,G,\dots,G^{S-1}$, so every later value repeats that combination of
earlier ones.
\end{proof}

\begin{theorem}
The recurrence \eqref{eq:conj} holds for every $n>15$, which is the statement of
Section~1.
\end{theorem}

\begin{proof}
Each $M_c$ was built directly from its defining condition. The vector $w$ was formed by
$15$ applications of $G$ in exact integer arithmetic and $\mathbf{v}^{\!\top}G^{m}w$
evaluated for $m=0,1,\dots$, again exactly. Every one of those integers vanishes from the
index corresponding to $n=15+1$ onwards, and the run continued until $S+1$ consecutive
zeros had been seen.
\end{proof}

\begin{remark}
The splitting is what makes the argument finite, and it is exact rather than an
over-count only because every array here has a subblock. For a one-column or one-row
shape there is no subblock at all, every value of $c$ would be vacuously admissible, and the
sum would count each array $12+1$ times; those shapes are excluded by $L\ge2$ and
$2\ge2$.
\end{remark}

\section{Verification}

Three checks, each independent of the algebra above.

First, the right-hand side of \eqref{eq:walk} was evaluated at every one of the $23$
published indices and agrees with the entry's DATA exactly, with no shift fitted. A double
count in the splitting, or a wrong bound on $c$, would show up here immediately.

Second, the recurrence \eqref{eq:conj} was evaluated on those published terms in exact
integer arithmetic, with no matrices involved, and vanishes wherever it applies.

Third, the annihilation test of Lemma~\ref{lem:crit} was carried to the full
Cayley--Hamilton bound in exact integer arithmetic rather than to a sampled prefix.

\begin{thebibliography}{9}
\bibitem{oeis} The OEIS Foundation, \emph{The On-Line Encyclopedia of Integer
Sequences}, \url{https://oeis.org}, 2026. Sequence A184021.
\bibitem{stanley} R.~P.~Stanley, \emph{Enumerative Combinatorics}, Volume 1, 2nd ed.,
Cambridge University Press, 2011. (Transfer-matrix method, Section 4.7.)
\bibitem{fs} P.~Flajolet and R.~Sedgewick, \emph{Analytic Combinatorics}, Cambridge
University Press, 2009.
\bibitem{horn} R.~A.~Horn and C.~R.~Johnson, \emph{Matrix Analysis}, 2nd ed., Cambridge
University Press, 2013.
\end{thebibliography}

\end{document}
